Here, we discuss prior work that conceptually and technically relates to ours. We then suggest natural directions for future work.

\paragraph{Multi-distribution learning.}
Many learning problems involve multiple data sources, for instance, when multiple agents generate their data independently. One can formulate these multi-distribution problems as standard learning/optimization problems by considering a mixture of their distributions, as in objective \ref{eq:utilitarian}. However, this approach often biases solutions toward dominant data sources, leading to poor performance on outliers—an issue stemming from statistical heterogeneity. This limitation motivates the study of multi-objective optimization problems~\citep{miettinen1999nonlinear, ehrgott2005multicriteria}, where each agent \( m \) has a distribution \( \mathcal{D}_m \) that defines its objective as \( \mathbb{E}_{z\sim \mathcal{D}_m}[f(\vx_m;z)] \), and where models \( \vx_m \) can vary across agents---a framework known as personalization. 

One of the earliest algorithms for such problems was introduced by \citet{blum2017collaborative}, where each agent's objective must be minimized to a pre-specified threshold \( \epsilon \) with high probability, framed within a PAC learning framework~\citep{valiant1984theory,vapnik2013nature}. Subsequent research has refined these algorithms, achieving optimal sample complexity guarantees for learning from multiple distributions~\citep{chen2018tight,nguyen2018improved,hanneke2019value,haghtalab2022demand,zhang2024optimal}. Our objectives \( \ref{eq:main_objective} \) and \( \ref{eq:interpolation} \) offer different approaches to multi-distribution learning, where data distributions correspond to empirical agent distributions. In particular, \citet{mohri2019agnostic} analyzed objective \( \ref{eq:main_objective} \) to establish generalization bounds for unknown mixtures of agents' distributions. 

Beyond sample efficiency, researchers have also examined other challenges, such as communication costs in large-scale distributed optimization~\citep{mcmahan2016s}. A particularly relevant study is that of \citet{bullins2021stochastic}, which employs an efficient distributed quadratic sub-solver~\citep{woodworth2020local,patel2024limits} to implement an inexact Newton method for optimizing quasi-self-concordant functions (see \Cref{defn:qsc}).

\paragraph{Group fairness.} 
Recently, interest in algorithmic fairness has intensified~\citep{barocas2016big, abebe2020roles, kasy2021fairness} with researchers exploring fairness across various domains, including supervised learning~\citep{calders2009building, dwork2012fairness, hardt2016equality, kusner2017counterfactual, goel2018non, ustun2019fairness}, resource allocation~\citep{bertsimas2011price, bertsimas2012efficiency, hooker2012combining, donahue2020fairness, manshadi2021fair}, scheduling~\citep{mulvany2021fair}, online matching~\citep{chierichetti2019matroids, ma2023fairness}, assortment planning~\citep{singh2018fairness, biega2018equity, singh2019policy, chen2022fair}, and facility location~\citep{gupta2022socially}. The extensive literature on algorithmic fairness falls into three main categories: (1) individual fairness, which ensures that similar individuals receive comparable predictions~\citep{dwork2012fairness, loi2019philosophical, chen2022fair}, (2) group fairness, which aims for equal treatment of different demographic groups, often in terms of resource allocation or performance parity~\citep{singh2018fairness, balseiro2021regularized}, and (3) subgroup fairness, which blends aspects of both individual and group fairness~\citep{kearns2018preventing, kearns2019empirical}. 

This paper focuses on a well-studied group fairness notion in machine learning literature: the group DRO problem~\citep{ben2013robust, duchi2016statistics, sagawa2019distributionally}. The idea of interpolating between robustness and utility is also common~\citep{golrezaei2024online} and closely related to multi-objective optimization, where scalarization~\citep{miettinen1999nonlinear, ehrgott2005multicriteria} helps recover desired solutions along the Pareto frontier. 

\paragraph{Linear programming and $\ell_p$ regression.} 
In the last several years, there has been a surge of work in obtaining second-order, condition-free algorithms for linear programming and $\ell_p$ regression \citep{bcll18,ls19,akps19,jls21}. Observe that $\ell_p$ regression is a special case of the problem we study in objective \eqref{eq:interpolation}, which is recovered when all $n_i=1$, and $\ell_{\infty}$ regression is captured by linear programming. Note that neither of these problem families is expressive enough to capture the objectives we study. In general, to achieve iteration complexities in the smaller of the two dimensions for these problems, it appears that a geometric understanding of the solution space is required --- these ideas were central to the improvements obtained by \citet{ls19,jls21} as well as our work.

\paragraph{Open problems.} 
Our work raises several open questions. One limitation of \Cref{mainthm:fair_regression_iteration_complexity} is that its iteration complexity is not high-accuracy, meaning its dependence on \( \eps \) is not \(\mathrm{polylog}(1/\eps)\). Designing a high-accuracy solver under the same conditions as \Cref{mainthm:gp_regression_iteration_complexity} with iteration complexity \(\widetilde{O}\inparen{\mathsf{poly}(\min\inbraces{\rank{\mA}, m}, \logv{\frac{1}{\eps}})}\) remains an open problem.

A more ambitious general goal is to design algorithms for convex quadratic programs with the aforementioned iteration complexity. This would generalize analogous results for linear programming \citep{ls19}. We view the current work as a first step towards this goal, as the objective \eqref{eq:main_objective} is a structured convex quadratic program for which we get an iteration complexity independent of $m$. It would also be interesting to consider other complexity measures beyond $\rank{\mA}$, for instance, assumptions about the ground-truth labeling vector $\vx_i^\star$ for each group's data $S_i$.

Finally, our results suggest that optimizing for ``$\ell_p$-interpolants'' between non-robust and robust objectives may be computationally easier than optimizing for the robust objective alone. A more precise statistical characterization of how robustness and utility trade-off as \( p \) varies in collaborative, fair, or multi-distributional learning settings would be valuable. Additionally, exploring interpolations or solution concepts along the Pareto frontier of the \( m \)-dimensional multi-objective optimization problem or other DRO notions (eg Wasserstein DRO \citep{bkm19, cpo20}) could yield further insights.
